\section{Open Questions and Next Work}
\label{sec:sources-wasm4pm-open-questions}

\subsection{Known Gaps}
\label{sec:sources-wasm4pm-gaps-known}

Ten structural gaps are documented in \texttt{GAP\_ANALYSIS.md} (dated 2026-05-31). The gaps
represent open research findings rather than implementation failures in \texttt{sources/wasm4pm}.

\textbf{GAP\_001} (Critical): The live \texttt{\textasciitilde/wasm4pm/} codebase carries zero
\texttt{wasm4pm-compat} dependencies. \texttt{Evidence<T,~State,~W>}, \texttt{Admission<T,~W>},
\texttt{Refusal<R,~W>}, and \texttt{Metric<KIND,NUM,DEN>} — all defined in \texttt{wasm4pm-compat}
— are absent from the engine side. The local \texttt{src/graduation.rs} bridges this seam for
the \texttt{sources/wasm4pm} codebase, but the relationship between the manufactured
\texttt{sources/wasm4pm} and the live \texttt{\textasciitilde/wasm4pm/} codebase is the
primary open research question.

\textbf{GAP\_002} (High): The live codebase's error surface is entirely \texttt{String}-typed.
Every \texttt{wasm4pm\_types::Error} variant carries a \texttt{String} payload with no named
structural law. \texttt{wasm4pm-compat} mandates named refusal reasons as first-class types.

\textbf{GAP\_003} / \textbf{GAP\_007} (High): The live \texttt{\textasciitilde/wasm4pm/} lacks
an Inductive Miner implementation (GAP\_003) and lacks alignment-based A* conformance
(GAP\_007). Both are present in \texttt{sources/wasm4pm} but absent in the live codebase —
reinforcing the two-codebase ambiguity of GAP\_001.

\textbf{GAP\_004} (High): No POWL discovery algorithm exists despite a spec-compliant
\texttt{Powl8Op} enum and \texttt{ChoiceGraph} type in the live codebase.

\textbf{GAP\_005} (High): No Object-Centric DFG discovery — all algorithms operate on flat
\texttt{EventLog} traces, not OCEL with \texttt{e2o}/\texttt{o2o} relations.

\textbf{GAP\_009} (High): BLAKE3 receipt fields in \texttt{Evidence} are initialized to
\texttt{[0u8;~32]}. The receipt infrastructure is structurally present but hash computation over
algorithm inputs is not wired.

Remaining gaps (\textbf{GAP\_006} Log Skeleton, \textbf{GAP\_008} OCPQ isolation,
\textbf{GAP\_010} ETConformance precision) are Medium severity and represent future algorithm
additions rather than architectural defects.

\subsection{Deferred Gates}
\label{sec:sources-wasm4pm-gaps-deferred}

The following proof gates are specified but not yet implemented. Variant enumeration API and
long-distance loop analysis (HIGH severity, Phase~2) are required to complete the Mining
Authority. Generalization metric computation and multi-model conformance for variant detection
(MEDIUM severity, Phase~3--5) are required to complete the Conformance Authority. Explicit
variable tracking in replay and batch replay optimization (MEDIUM/LOW severity, Phase~5--6) are
required to complete the Replay Authority. Model versioning/rollback and concurrent model
deployment for A/B testing (MEDIUM severity, Phase~5) are required to complete the Lifecycle
Authority. Thread-local receipt accumulation for concurrent replay (no severity assigned in the
research verdict) and archival recovery via board override signature are deferred to Phase~5.

\subsection{Research Risks}
\label{sec:sources-wasm4pm-gaps-risks}

The primary research risk is the \emph{two-codebase ambiguity}: whether \texttt{sources/wasm4pm}
(the manufactured codebase with full \texttt{wasm4pm-compat} integration) or the live
\texttt{\textasciitilde/wasm4pm/} workspace (the engineering codebase without compat consumption)
constitutes the authoritative wasm4pm. The dissertation must resolve this question before making
final claims about the execution authority. If the live codebase is authoritative, then GAP\_001
through GAP\_010 are live defects, not research findings in a separate artifact.

A secondary risk is the \emph{receipt-field zero-initialization}: BLAKE3 hash fields initialized
to \texttt{[0u8;~32]} mean that receipts are currently structurally present but cryptographically
empty. If a spot-audit is conducted before the hash computation is wired, the audit will confirm
structural correctness but cannot confirm cryptographic binding. This must be resolved before any
capital-flow claim is made on the basis of receipt chain integrity.

A third risk is \emph{compile-time vs. runtime metric gap}: \texttt{Between01<NUM,DEN>} works
only for const-generic parameters, while all runtime algorithm outputs require the separate
\texttt{RuntimeBetween01} path. The bridge between the two is not systematic across the
\texttt{Evidence} wrapper.

\subsection{Next Receipted Work}
\label{sec:sources-wasm4pm-gaps-next-work}

The following work items require receipted completion before the dissertation can close the
open questions. (1) Resolve the two-codebase ambiguity: determine whether \texttt{sources/wasm4pm}
is a specification artifact or the authoritative execution engine, and document the resolution in
a new checkpoint file. (2) Wire BLAKE3 hash computation over algorithm inputs in all
\texttt{Evidence<T,~S,~W>} constructors, replacing \texttt{[0u8;~32]} initialization with
actual hash computation; seal the result with a new manufacturing receipt. (3) Implement
\texttt{wasm4pm-compat} consumption in the live \texttt{\textasciitilde/wasm4pm/} workspace
(closing GAP\_001); issue a checkpoint receipt upon successful compilation with the
\texttt{GraduateToWasm4pm} trait implemented. (4) Implement the Inductive Miner cut-detection
algorithm in the live codebase (closing GAP\_003); issue a discovery receipt confirming
block-structured soundness by construction. (5) Wire the A* alignment engine into the live
conformance surface (closing GAP\_007); issue a conformance receipt confirming the fitness
equation is computed from actual alignment moves, not from token replay. All five items must
be completed and receipted before any ALIVE verdict is issued for the live
\texttt{\textasciitilde/wasm4pm/} workspace. [FUTURE WORK]
